\clearpage
\chapter*{강연 190의 정오표}
\addcontentsline{toc}{chapter}{강연 190의 정오표}
\setcounter{footnote}{0}
\src{300}

\noindent\textbf{3. 대수기하학에서의 하강 기법과 존재 정리
[TDTE I]. I: 일반론. 충실 평탄 사상에 의한 하강.}

\begin{center}
[부르바키 세미나, 제12권, 1959/60, 제190호, 29쪽.]
\end{center}

\noindent\textbf{1쪽, 10행.} --- 이제 하강 기법이 ``대수기하학의
존재 정리 대부분의 기초''라고 말하는 것은 지나쳐 보인다. 이는 TDTE
I--VI 계열의 처음 두 강연에서 다루는 비사영적 기법들에는 상당한
정도로 타당하지만, 사영적 기법들(TDTE IV, V, VI)에는 그렇지 않다.

\noindent\textbf{5쪽, 16행.} --- $\alpha$가 $\mathcal F$-하강 사상이라고
가정할 필요는 없다.

\noindent\textbf{20쪽, 비고.} --- 준유한이고 에탈이며 전사인 사상
$S'\longrightarrow S$, 또는
$\operatorname{Spec}(\widehat A)\longrightarrow\operatorname{Spec}(A)$
꼴의 사상은 언제나 엄밀 하강 사상인 것은 아니다. 이는 $A$가
대수적으로 닫힌 체 $k$ 위의 대수곡선의 국소환이고
$S=\operatorname{Spec}(A)$인 경우에도 마찬가지이다. 따라서 $X$가 밑이
$S$이고 타원곡선 $E$를 구조군으로 하는 주다발이 되게 하는 고유하고
매끄러운 사상 $f:X\longrightarrow S$로서,
$f':X'\longrightarrow S'$은 사영이지만 $f$는 사영이 아닌 것을 찾을
수 있다. 그러므로 이는 아벨 스킴을 구조군으로 하는 등상적이지 않은
동차 주다발의 예이기도 하다.

\noindent\textbf{26쪽, 9행.} --- ``Igusa--Lang'' 대신 ``Chow--Lang''으로
읽어야 한다.

하강 이론의 여러 세부사항에 대해서는 SGA VI, VII, VIII를 보라.

% FGA Canon print-preservation apparatus: atomic integration R1.
\par\smallskip
\begingroup\footnotesize
\noindent\textit{FGA-CANON-ERR-0031. 인쇄본의 독법은 C094에 계속 기록되어 있으며, FGA-CANON-DISP-0033이 현행 본문의 최소 수정을 규율한다.}\par
\endgroup
